\documentclass[11pt]{article}
\usepackage[utf8]{inputenc}
\usepackage[margin=1in]{geometry}
\usepackage{graphicx}
\usepackage{booktabs}
\usepackage{xcolor}
\usepackage{hyperref}
\usepackage{listings}
\usepackage{amsmath}
\usepackage{float}

\definecolor{codegray}{gray}{0.95}
\lstset{
    backgroundcolor=\color{codegray},
    basicstyle=\ttfamily\footnotesize,
    breaklines=true,
    frame=single,
    columns=fullflexible
}

\title{U1T01: Real-Time Analytics with PostgreSQL CDC \& ClickHouse}
\date{}

\begin{document}

\maketitle
\thispagestyle{empty}

\vspace{-1.5cm}

\begin{center}

\begin{minipage}[t]{0.30\textwidth}
\centering
\textbf{Student} \\[0.2em]
Mauricio Gabriel Ram\'{i}rez Rubio \\
Data Engineering --- 8th Term \\
Universidad Polit\'{e}cnica de Yucat\'{a}n \\
Km.\ 4.5 Carretera M\'{e}rida --- Tetiz \\
Tablaje Catastral 4448. CP 97357 \\
Uc\'{u}, Yucat\'{a}n. M\'{e}xico \\
Email: 2309192@upy.edu.mx
\end{minipage}
\hfill
\begin{minipage}[t]{0.30\textwidth}
\centering
\textbf{Student} \\[0.2em]
Cesar Arturo Balam Euan \\
Data Engineering --- 8th Term \\
Universidad Polit\'{e}cnica de Yucat\'{a}n \\
Km.\ 4.5 Carretera M\'{e}rida --- Tetiz \\
Tablaje Catastral 4448. CP 97357 \\
Uc\'{u}, Yucat\'{a}n. M\'{e}xico \\
Email: 2309017@upy.edu.mx
\end{minipage}
\hfill
\begin{minipage}[t]{0.30\textwidth}
\centering
\textbf{Student} \\[0.2em]
Rodrigo Garc\'{i}a Mart\'{i}nez \\
Data Engineering --- 8th Term \\
Universidad Polit\'{e}cnica de Yucat\'{a}n \\
Km.\ 4.5 Carretera M\'{e}rida --- Tetiz \\
Tablaje Catastral 4448. CP 97357 \\
Uc\'{u}, Yucat\'{a}n. M\'{e}xico \\
Email: 2309091@upy.edu.mx
\end{minipage}

\vspace{1.2cm}

\begin{minipage}[t]{0.30\textwidth}
\centering
\textbf{Student} \\[0.2em]
Diego De Gante P\'{e}rez \\
Data Engineering --- 8th Term \\
Universidad Polit\'{e}cnica de Yucat\'{a}n \\
Km.\ 4.5 Carretera M\'{e}rida --- Tetiz \\
Tablaje Catastral 4448. CP 97357 \\
Uc\'{u}, Yucat\'{a}n. M\'{e}xico \\
Email: 2309067@upy.edu.mx
\end{minipage}
\hfill
\begin{minipage}[t]{0.30\textwidth}
\centering
\textbf{Student} \\[0.2em]
Roselyn Guadalupe Polanco Gonz\'{a}lez \\
Data Engineering --- 8th Term \\
Universidad Polit\'{e}cnica de Yucat\'{a}n \\
Km.\ 4.5 Carretera M\'{e}rida --- Tetiz \\
Tablaje Catastral 4448. CP 97357 \\
Uc\'{u}, Yucat\'{a}n. M\'{e}xico \\
Email: 2309184@upy.edu.mx
\end{minipage}
\hfill
\begin{minipage}[t]{0.30\textwidth}
\centering
\textbf{Professor} \\[0.2em]
G\'{o}mez Dexter \\
Trends in Data Science \\
Universidad Polit\'{e}cnica de Yucat\'{a}n \\
Km.\ 4.5 Carretera M\'{e}rida --- Tetiz \\
Tablaje Catastral 4448. CP 97357 \\
Uc\'{u}, Yucat\'{a}n. M\'{e}xico \\
Email: gomez.dexter@upy.edu.mx
\end{minipage}

\end{center}

\vspace{1cm}

\section{Introduction}

Modern data architectures separate transactional (OLTP) and analytical (OLAP)
workloads across two different database engines, since no single engine
covers both scenarios efficiently. This project implements Change Data
Capture (CDC) between a PostgreSQL OLTP database and a ClickHouse OLAP
database, using a custom, dependency-light Python service that reads
PostgreSQL's Write-Ahead Log (WAL) directly through logical decoding,
instead of relying on an external orchestration platform. This mechanism
is the same underlying primitive that tools such as PeerDB or Debezium use
internally.

\section{Part 1 --- Relational Modeling}

An e-commerce domain was modeled in Third Normal Form (3NF) with six
tables and clear primary/foreign key relationships:

\begin{itemize}
    \item \textbf{categories} --- product category catalog
    \item \textbf{customers} --- registered users
    \item \textbf{products} --- catalog items, each belonging to one category
    \item \textbf{orders} --- purchase headers, referencing a customer
    \item \textbf{order\_items} --- line items, resolving the orders$\leftrightarrow$products
          many-to-many relationship
    \item \textbf{payments} --- one payment record per order
\end{itemize}

\textbf{Normalization rationale:} \texttt{category\_id} is stored on
\texttt{products} rather than the category name, avoiding a transitive
dependency. \texttt{unit\_price} is captured on \texttt{order\_items} at
the time of purchase, rather than being derived from the live
\texttt{products.price} value. Payment attributes (\texttt{method},
\texttt{status}) live in their own table rather than being repeated on
\texttt{orders}.

\subsection*{ER Diagram}
\begin{figure}[H]
    \centering
    \fbox{\parbox{0.9\textwidth}{\centering \vspace{6cm}
    [ INSERT ER DIAGRAM SCREENSHOT HERE --- e.g. from dbdiagram.io,
    pgAdmin's ER tool, or a draw.io export of the 6-table schema ]
    \vspace{6cm}}}
    \caption{Entity-Relationship diagram of the e-commerce schema (3NF).}
\end{figure}

\subsection*{Data Generator}
A continuous Python generator (\texttt{generator.py}) produces a weighted
mix of operations against PostgreSQL: 35\% new orders (INSERT across
\texttt{orders}, \texttt{order\_items}, \texttt{payments}), 25\% order
status advances (UPDATE), 15\% product price/stock changes (UPDATE), 10\%
new customer registrations (INSERT), 10\% payment status transitions
(UPDATE), and 5\% cleanup of old cancelled orders (DELETE, cascading to
\texttt{order\_items}). The generator ran continuously and produced well
over 300{,}000 operations during testing.

\begin{figure}[H]
    \centering
    \fbox{\parbox{0.9\textwidth}{\centering \vspace{4cm}
    [ INSERT SCREENSHOT: \texttt{docker logs ecommerce\_generator --tail 20}
    showing the ops counter and rate, e.g. ``ops=350{,}000 rate=78 ops/s'' ]
    \vspace{4cm}}}
    \caption{Data generator running continuously.}
\end{figure}

\section{Part 2 --- CDC Pipeline Setup}

\subsection*{Why a custom CDC worker instead of an external platform}
PeerDB was initially attempted, but its architecture (Nexus query layer,
Flow API, Flow Workers, a Temporal orchestration engine, a separate
catalog database, and MinIO for staging) introduced multiple points of
failure in a local Docker Desktop environment --- specifically, workflows
never reached the Temporal task queue due to a gRPC hand-off issue between
\texttt{peerdb-server} and \texttt{flow\_api}. The assignment explicitly
allows using ``any other open-source tool available,'' so the CDC pipeline
was implemented directly against PostgreSQL's own logical decoding
mechanism --- the same primitive PeerDB and Debezium build on top of.

\subsection*{PostgreSQL configuration}
\begin{enumerate}
    \item \texttt{wal\_level=logical} enables logical decoding.
    \item \texttt{max\_replication\_slots} / \texttt{max\_wal\_senders} sized for the workload.
    \item The \texttt{wal2json} output plugin was installed via the official
          \texttt{postgresql-15-wal2json} apt package on top of the standard
          \texttt{postgres:15} image.
    \item As of the running PostgreSQL version, output plugins must be
          explicitly allow-listed via \texttt{output\_plugin\_libraries=wal2json}
          (a security hardening measure) before a replication slot can use them.
    \item \texttt{REPLICA IDENTITY FULL} was set on all six tables so that
          UPDATE/DELETE events carry the complete old row --- required to
          build a full ``tombstone'' row in ClickHouse when a DELETE occurs.
\end{enumerate}

\subsection*{The CDC worker (\texttt{cdc\_worker.py})}
A single Python service performs the entire pipeline:
\begin{enumerate}
    \item Opens a logical replication slot with the \texttt{wal2json} plugin
          using \texttt{psycopg2}'s replication protocol support.
    \item On first run, performs a full initial copy of all six tables into
          ClickHouse with version \texttt{1}.
    \item Streams every subsequent WAL change and decodes the JSON payload
          wal2json emits (\texttt{insert}/\texttt{update}/\texttt{delete}
          events with column names and values).
    \item Buffers changes per table and flushes them to ClickHouse (via the
          HTTP interface, \texttt{clickhouse-connect}) at least every 2
          seconds, or sooner if a buffer fills up.
    \item Each row is written with a monotonically increasing
          \texttt{\_version} (nanosecond timestamp) and an
          \texttt{\_is\_deleted} flag, so ClickHouse's
          \texttt{ReplacingMergeTree} engine keeps only the latest version
          of each primary key when queried with \texttt{FINAL}.
\end{enumerate}

\begin{figure}[H]
    \centering
    \fbox{\parbox{0.9\textwidth}{\centering \vspace{5cm}
    [ INSERT SCREENSHOT: \texttt{docker logs cdc\_worker} showing
    ``Created new replication slot'', ``Starting initial full load of all
    tables'', and repeated ``Flushed N rows -> ecommerce.<table>'' lines ]
    \vspace{5cm}}}
    \caption{CDC worker performing the initial load and streaming changes.}
\end{figure}

\section{Part 3 --- Analytical Layer \& CDC Verification}

ClickHouse tables mirror the PostgreSQL schema, using the
\texttt{ReplacingMergeTree} engine keyed on each table's primary key, with
two extra columns: \texttt{\_version} (drives deduplication) and
\texttt{\_is\_deleted} (marks tombstones from DELETE events). Since
ClickHouse does not support in-place updates or deletes the way an OLTP
engine does, every query must either use the \texttt{FINAL} modifier or
filter \texttt{\_is\_deleted = 0} to see the current state of a row.

\subsection*{Verifying propagation}
An \texttt{UPDATE} was issued directly against PostgreSQL and confirmed to
appear in ClickHouse within the worker's flush interval ($\leq 2$ seconds):

\begin{lstlisting}
-- PostgreSQL
UPDATE products SET price = 999.99 WHERE product_id = 1;

-- ClickHouse (queried ~3s later)
SELECT product_id, price FROM ecommerce.products FINAL WHERE product_id = 1;
\end{lstlisting}

\begin{figure}[H]
    \centering
    \fbox{\parbox{0.9\textwidth}{\centering \vspace{4cm}
    [ INSERT SCREENSHOT: terminal showing the UPDATE in Postgres and the
    matching SELECT ... FINAL result in ClickHouse with the new price ]
    \vspace{4cm}}}
    \caption{An UPDATE issued in PostgreSQL reflected in ClickHouse.}
\end{figure}

\subsection*{Row count comparison}
Row counts between the two systems stay closely in sync while the
generator is actively producing data (a small, expected lag exists because
counts are two snapshots taken seconds apart while $\sim$80 ops/s are being
written):

\begin{figure}[H]
    \centering
    \fbox{\parbox{0.9\textwidth}{\centering \vspace{3cm}
    [ INSERT SCREENSHOT: \texttt{SELECT COUNT(*) FROM orders;} on both
    PostgreSQL and ClickHouse showing close row counts ]
    \vspace{3cm}}}
    \caption{Order counts staying in sync between OLTP and OLAP.}
\end{figure}

\section{Part 4 --- Benchmarking}

Five analytical queries involving joins, aggregation, filtering, and
window functions were run against both systems while the generator was
actively writing data.

\begin{enumerate}
    \item \textbf{Q1:} Daily revenue by product category, last 30 days (4-way join).
    \item \textbf{Q2:} Customer lifetime value, top 50 spenders (3-way join, HAVING).
    \item \textbf{Q3:} Order fulfillment rate and average ticket by city (conditional aggregation).
    \item \textbf{Q4:} Product performance --- revenue and units sold (4-way join).
    \item \textbf{Q5:} Payment method breakdown with window function (share \% per method).
\end{enumerate}

\subsection*{Results}

\begin{table}[H]
\centering
\caption{Execution time comparison: PostgreSQL (\texttt{EXPLAIN ANALYZE}) vs.\ ClickHouse.}
\begin{tabular}{lrrl}
\toprule
\textbf{Query} & \textbf{PostgreSQL (ms)} & \textbf{ClickHouse (ms)} & \textbf{Notes} \\
\midrule
Q1 --- Daily revenue by category   & 329.6  & 427.4  & 4-way join, $\sim$82k rows scanned \\
Q2 --- Customer lifetime value     & 301.5  & --     & 3-way join, HAVING \\
Q3 --- Fulfillment rate by city    & 54.9   & 482.0  & simple aggregation \\
Q4 --- Product performance         & 612.8  & --     & heaviest query, 4-way join \\
Q5 --- Payment method breakdown    & 20.3   & 238.6  & window function \\
\bottomrule
\end{tabular}
\end{table}

\textit{Note: ClickHouse timings were captured with PowerShell's
\texttt{Measure-Command} wrapping \texttt{docker exec}, which includes
container-invocation overhead not present in PostgreSQL's
\texttt{EXPLAIN ANALYZE} figure; Q2 and Q4 were not timed this way due to
time constraints and could be added by running the same
\texttt{Measure-Command} pattern against those two queries.}

\begin{figure}[H]
    \centering
    \fbox{\parbox{0.9\textwidth}{\centering \vspace{4cm}
    [ INSERT SCREENSHOT: PostgreSQL EXPLAIN ANALYZE output, e.g. Q4 showing
    the Merge Join / Incremental Sort / Seq Scan plan ]
    \vspace{4cm}}}
    \caption{PostgreSQL query plan example (Q4 --- product performance).}
\end{figure}

\begin{figure}[H]
    \centering
    \fbox{\parbox{0.9\textwidth}{\centering \vspace{4cm}
    [ INSERT SCREENSHOT: ClickHouse EXPLAIN output showing
    ReadFromMergeTree + Join (JOIN FillRightFirst) steps ]
    \vspace{4cm}}}
    \caption{ClickHouse query plan example (same query).}
\end{figure}

\subsection*{Why the plans look different}

\textbf{PostgreSQL} is a row-oriented engine optimized for OLTP: every
query plan inspected relies on \texttt{Seq Scan} (sequential, row-by-row
scans) combined with \texttt{Hash Join} and \texttt{Sort} nodes. Because
PostgreSQL has no columnar storage, it must read entire rows even when a
query only needs a few columns, and sorting large intermediate result sets
(seen as ``Sort Method: external merge, Disk: \dots'' in Q1 and Q2) spills
to disk when it exceeds \texttt{work\_mem}.

\textbf{ClickHouse} is a column-oriented engine built for OLAP: every plan
uses \texttt{ReadFromMergeTree}, which reads only the columns referenced
by the query directly off disk in a vectorized, batch-oriented fashion.
Its own \texttt{Join (JOIN FillRightFirst)} operator builds an in-memory
hash table from the right-hand relation and streams the left-hand
relation through it, avoiding the need to materialize full rows before
filtering. This columnar, vectorized execution model is why ClickHouse
scales far better on large-scale aggregation over a small number of
columns, whereas PostgreSQL is architecturally suited for point lookups,
transactional consistency, and workloads that touch full rows one at a
time --- exactly the OLTP vs.\ OLAP split this assignment set out to
demonstrate empirically.

\section{Conclusion}

This project implemented a full CDC pipeline between PostgreSQL and
ClickHouse using PostgreSQL's native logical decoding (\texttt{wal2json})
rather than an external orchestration platform, verified that inserts,
updates, and deletes propagate correctly and promptly, and benchmarked
five analytical queries on both systems. The results confirm the expected
architectural trade-off: PostgreSQL is well suited to the transactional
workload it was serving, while ClickHouse's columnar, vectorized engine is
built for the kind of aggregation-heavy analytical queries this benchmark
exercised.

\section*{Repository}
Full source code, Docker Compose configuration, and SQL scripts are
available at: \url{[ INSERT YOUR GITHUB REPO LINK HERE ]}

\end{document}
